\section{Uncertainty quantification}
\label{sec:systematics}

The uncertainty budget combines three independent blocks --- systematic
(universe), statistical (bootstrap), and ML (seedscan) --- summed assuming
independence (distinct RNGs / physics sources). This section describes each
block and their combination for the 2D measurement; the 3D and 4D campaigns
(\S\ref{sec:3d-syst}, \S\ref{sec:4d}) repeat the identical procedure on the
higher-dimensional grids. Full mathematical procedures, with worked examples
on the reported numbers, are in App.~\ref{app:statmethods}. The headline is
that our \emph{standalone} combined budget matches the published total.

\begin{table}[H]
  \centering
  \begin{tabular}{@{}lccc@{}}
    \toprule
    Component & $N$ & $\sqrt{\mathrm{Tr}\,C}$ & per-bin median rel.\ $\sigma$ \\
    \midrule
    ML noise (lgbm seedscan) & 10 & \num{5.06e-41} & \SI{0.166}{\percent} \\
    Statistical (Poisson bootstrap) & 300 & \num{1.82e-40} & \SI{0.549}{\percent} \\
    Systematic (universe sweep $+$ \SI{1.4}{\percent} norm) & $187{+}1$ & \num{3.21e-39} & \SI{6.830}{\percent} \\
    \textbf{Combined (block sum)} & --- & \num{3.22e-39} & \SI{6.865}{\percent} \\
    Paper TotalCov (reference) & --- & \num{2.68e-39} & \SI{6.86}{\percent} \\
    \bottomrule
  \end{tabular}
  \caption{Stage-2 UQ envelope on the 205 reported bins. The standalone combined
  budget (\SI{6.865}{\percent}) matches the published total (\SI{6.86}{\percent}).}
  \label{tab:uq}
\end{table}

\subsection{Systematic universes}
\label{sec:uq-universes}
The detector and model systematics use the MAT-MINERvA universe machinery: 44
error bands (6 lateral --- shifted reconstructed kinematics --- and 38
vertical --- per-event reweights), 187 universes in total, identical to the
published analysis's \texttt{GetStandardSystematics} set. Each universe is a
\emph{full re-unfold}: the event loop dumps the universe's weights and (for
lateral bands) shifted observables, and the entire OmniFold chain runs on
them with the central value's seeds and configuration, so the systematic
propagates through the learned reweighting itself rather than through a frozen
response. Pair ($\pm1\sigma$) bands use the MINERvA sum-of-squares convention
$\tfrac12(\delta_+\delta_+^{\!\top}+\delta_-\delta_-^{\!\top}$); multi-universe
bands (the 100-universe PPFX flux multisim) use the sample covariance. The
construction is MAT-conformant (\texttt{MnvVertErrorBand::CalcCovMx}; verified
byte-for-byte against MAT \texttt{MnvH1D::GetTotalErrorMatrix}, max rel.\ diff
$\sim10^{-17}$). Top bands: Flux \SI{\fluxBand}{\percent}, Muon\_Energy\_MINOS
\SI{2.31}{\percent}, MinosEfficiency \SI{1.47}{\percent},
Muon\_Energy\_MINERvA \SI{1.29}{\percent}. Figure~\ref{fig:uqbands} shows the
grouped fractional-uncertainty projections in the published Fig.~6/7 style
(the full $205\times205$ covariance projected exactly, $C_{1D}=PC_{2D}P^{\top}$,
before dividing by the 1D central value).

\begin{figure}[H]
  \centering
  \includegraphics[width=0.49\textwidth]{MEFHC_fig6_7_uncertainty_pt.png}
  \includegraphics[width=0.49\textwidth]{MEFHC_fig6_7_uncertainty_pz.png}
  \caption{Fractional uncertainty on the 2D cross section projected onto $\pT$
  (left) and $\pPar$ (right), grouped by source in the style of Figs.~6--7 of
  Ref.~\cite{MINERvA:2021owq}: flux, interaction models, normalisation,
  hadronic response, muon reconstruction, and the statistical band.}
  \label{fig:uqbands}
\end{figure}

\paragraph{Flux $1/\Phi$ normalisation.}
The dominant flux term enters as $\sigma\propto1/\Phi$; varying $\Phi$ per flux
universe (rather than dividing every universe by the CV $\Phi$) raised the Flux
band from $\sim\SI{1}{\percent}$ to \SI{\fluxBand}{\percent} and brought the standalone
combined budget to \SI{\uqMedian}{\percent}. The \SI{1.4}{\percent} target-nucleon
normalisation is a fully-correlated rank-1 band. The 100-universe flux
covariance has a participation-ratio effective rank of $\approx1$ (one coherent
normalisation mode carries \SI{99.6}{\percent} of the variance), and its leading
mode correlates at $\rho=\rhoFluxMuon$ with the rank-1 muon-energy-scale band:
the rederived Bashyal~\cite{Bashyal:2021tqu} joint block, whose cross term is
rank-2, with two independent estimators giving $\rho=0.87$--$0.90$. Adding the
block leaves the systematic rank unchanged --- its directions already lie within
the span of the universe covariance (App.~\ref{sec:rank}).

\subsection{Statistical uncertainty: Poisson bootstrap}
\label{sec:uq-boot}
The statistical block is a 300-replica per-event Poisson(1) weight bootstrap
applied jointly to data and simulation, with two independent sub-RNGs (data
and MC are independent variance sources), the MC reco and truth weights riding
a \emph{single} per-event draw (preserving the within-event migration
correlation), and the GBDT seed pinned per replica so the bootstrap and
seedscan variances are separable. Each replica is a full re-unfold; the
covariance is the sample covariance over replicas
(Fig.~\ref{fig:bootstrap}). Our bootstrap block is genuinely smaller than the
published statistical block ($\approx2.5\times$ by trace) --- an unbinned
statistical efficiency confirmed by the data-only stream of the data/MC split
bootstrap ($C_{\mathrm{data}}/\texttt{StatOnly}=0.356$), not a missing term
(App.~\ref{sec:rank}).

\begin{figure}[H]
  \centering
  \includegraphics[width=0.49\textwidth]{uq_spread_2d.png}
  \includegraphics[width=0.49\textwidth]{uq_corr_2d.png}
  \caption{300-replica Poisson bootstrap: per-bin relative spread of the
  unfolded 2D cross section (left) and the bootstrap correlation matrix on the
  205 reported bins (right).}
  \label{fig:bootstrap}
\end{figure}

\subsection{ML/optimisation uncertainty: seedscan}
\label{sec:uq-ml}
ML stochasticity is budgeted by re-unfolding with 10 independent estimator
seeds at fixed data and fixed universe (Fig.~\ref{fig:seedscan}); the sample
covariance over seeds is the ML block. It is the smallest block
(median \SI{0.166}{\percent} per bin). Two refinements harden it: (i) a
\emph{train/test-split} seedscan, which also re-draws the internal training
split, gives a band $1.24\times$ the pure-seed band, moving the combined 3D
budget by only $+0.04\%$ in $\sqrt{\mathrm{Tr}}$ --- the published budgets use
the split-inclusive band; (ii) the ensemble-mean central value (averaging the
seed ensemble) agrees with the production single-run CV to \SI{0.28}{\percent}.

\begin{figure}[H]
  \centering
  \includegraphics[width=0.49\textwidth]{seedscan_spread_2d.png}
  \includegraphics[width=0.49\textwidth]{seedscan_band_pt.png}
  \caption{ML seedscan: per-bin relative spread over 10 estimator seeds (left)
  and the resulting fractional band projected onto $\pT$ (right).}
  \label{fig:seedscan}
\end{figure}

\subsection{Combination and the unified-throw test}
\label{sec:uq-combine}
The published covariance is the block sum
$C = C_{\mathrm{syst}} + C_{\mathrm{stat}} + C_{\mathrm{ML}}$, which assumes
the blocks are independent \emph{and} that the unfolding responds linearly to
simultaneous perturbations (so that per-band covariances add). The second
assumption is tested directly by the \emph{unified throw}
(\S\ref{sec:4d}): joint re-unfolds with all reweightable knob bands and the
flux multisim perturbed simultaneously per event. In 2D/3D the block sum is
consistent with the throws; on the 4D grid it \emph{underestimates} the
vertical systematic by a factor $\sim2$, concentrated in the high-$\pT$,
lowest-$\Eavail$ corner, and the published 4D covariance therefore adopts the
unified-throw magnitude via a PSD-safe per-bin inflation
(\S\ref{sec:4d}). This is a direct, quantitative validation (and, where
broken, correction) of the linearity assumption underlying the standard
block-sum construction.

\paragraph{$\chi^{2}$ caveat.}
The combined paper$+$ours $\chi^{2}/\mathrm{ndf}=1.481$ double-counts systematics
shared with the paper (flux/GENIE/MINOS); the ours-only inverse-covariance
$\chi^{2}$ is evaluated only with a truncated-spectral convention on the
well-constrained subspace. We therefore quote per-bin pulls, the regularisation
scan, and the truncated result together rather than treating an unconstrained
pseudo-inverse as a standalone goodness-of-fit (App.~\ref{sec:rank}).
